%
% music-playlist.tex
%
% Z Specification for the Vox self-driving music playlist (bead vox-1rxb)
% Formal model of the eager-fill / auto-advance / rotate-at-12 state machine
%

\documentclass[a4paper,10pt,fleqn]{article}
\usepackage[margin=1in]{geometry}
\usepackage{fuzz}
\usepackage[colorlinks=true,linkcolor=blue,citecolor=blue,urlcolor=blue]{hyperref}

\begin{document}

\title{Vox Self-Driving Music Playlist: A Z Specification}
\author{Formal Model of vox-1rxb (eager fill, auto-advance, rotate-at-12)}
\date{July 2026}
\hypersetup{
  pdftitle={Vox Self-Driving Music Playlist: A Z Specification},
  pdfauthor={Punt Labs},
  pdfsubject={Z Specification},
  pdfcreator={fuzz/probcli}
}
\maketitle

\tableofcontents
\newpage

\section{Introduction}

This specification pins the target behaviour of the vox background-music
playlist for a single \emph{(vibe, style)} pool. On turn-on the daemon
generates the first track and starts playing it, then generates the
remaining tracks in the background (one at a time) up to a fixed pool
size. Playback advances automatically as each track ends; once the pool
is full, playback rotates among the saved tracks with no further
generation. The model exists to make the transitions and invariants
unambiguous before implementation, and to expose any underspecified or
contradictory cases (see Section~\ref{sec:findings}).

\section{Basic Types}

A track is an opaque saved-audio identifier (in the implementation, the
content-addressed filename within a pool). A session identifies the agent
that turned music on -- the \emph{owner} whose forwarded vibe and skip
commands the daemon honours.

\begin{zed}
  [TRACK, SESSION]
\end{zed}

\section{Free Types}

The playback mode is a four-state machine.

\begin{zed}
  Mode ::= off | generatingFirst | playingFilling | playingRotating
\end{zed}

A Z boolean (avoids the ProB/B keyword clash with \texttt{BOOL}).

\begin{zed}
  ZBOOL ::= ztrue | zfalse
\end{zed}

\section{Global Constants}

The pool fills to \textbf{maxPool} distinct tracks, after which
generation stops and playback rotates. The design value is 12.

\begin{axdef}
  maxPool : \nat
  \where
  maxPool = 12
\end{axdef}

\section{State}

\texttt{playing} and \texttt{lastPlayed} are modelled as sets of at most
one track (an ``optional track''): empty when nothing is playing / no
track has yet been played.

\texttt{owner} is the session that turned music on, modelled as a set of at
most one session: empty exactly when music is off. It gates who may change
the vibe or skip.

\begin{schema}{Playlist}
  pool : \power TRACK \\
  playing : \power TRACK \\
  lastPlayed : \power TRACK \\
  owner : \power SESSION \\
  mode : Mode \\
  filling : ZBOOL
  \where
  \# pool \leq maxPool \\
  \# playing \leq 1 \\
  \# lastPlayed \leq 1 \\
  \# owner \leq 1 \\
  playing \subseteq pool \\
  lastPlayed \subseteq pool \\
  owner = \emptyset \iff mode = off \\
  mode = off \implies (playing = \emptyset \land filling = zfalse) \\
  mode = generatingFirst \implies \\
  \quad~ (pool = \emptyset \land playing = \emptyset \land filling = ztrue) \\
  mode = playingFilling \implies \\
  \quad~ (\# playing = 1 \land \# pool \geq 1 \\
  \quad~ \land \# pool < maxPool \land filling = ztrue) \\
  mode = playingRotating \implies \\
  \quad~ (\# playing = 1 \land \# pool = maxPool \land filling = zfalse) \\
  filling = ztrue \implies \\
  \quad~ (mode = generatingFirst \lor mode = playingFilling)
\end{schema}

\section{Initialisation}

The daemon starts idle with an empty pool. \texttt{InitFromDisk}
models a restart where a pool of already-saved tracks is loaded (still
idle until turned on).

\begin{schema}{Init}
  Playlist'
  \where
  pool' = \emptyset \\
  playing' = \emptyset \\
  lastPlayed' = \emptyset \\
  owner' = \emptyset \\
  mode' = off \\
  filling' = zfalse
\end{schema}

\begin{schema}{InitFromDisk}
  Playlist' \\
  diskPool? : \power TRACK
  \where
  \# diskPool? \leq maxPool \\
  pool' = diskPool? \\
  playing' = \emptyset \\
  lastPlayed' = \emptyset \\
  owner' = \emptyset \\
  mode' = off \\
  filling' = zfalse
\end{schema}

\section{Operations}

\subsection{TurnOn}

From \texttt{off}, enter the pool. An empty pool begins generating the
first track; a partial pool starts playing and resumes the background
fill; a full pool starts playing and rotates. \texttt{who?} becomes the
owner. (Restart is \texttt{InitFromDisk} followed by \texttt{TurnOn}.)

\begin{schema}{TurnOn}
  \Delta Playlist \\
  who? : SESSION
  \where
  mode = off \\
  pool' = pool \\
  lastPlayed' = \emptyset \\
  owner' = \{ who? \} \\
  (pool = \emptyset \implies \\
    \quad~ (mode' = generatingFirst \land playing' = \emptyset \\
  \quad~ \land filling' = ztrue)) \\
  (\# pool \geq 1 \land \# pool < maxPool \implies \\
    \quad~ (mode' = playingFilling \land playing' \subseteq pool \\
  \quad~ \land \# playing' = 1 \land filling' = ztrue)) \\
  (\# pool = maxPool \implies \\
    \quad~ (mode' = playingRotating \land playing' \subseteq pool \\
  \quad~ \land \# playing' = 1 \land filling' = zfalse))
\end{schema}

\subsection{FillProduces}

The background fill delivers one fresh track. It is enabled only while
\texttt{filling} and the pool is not yet full, and the track must be new
(the deterministic collision-free naming from bas7 guarantees this). The
first delivered track starts playback; reaching \texttt{maxPool} stops
the fill and switches to rotation.

\begin{schema}{FillProduces}
  \Delta Playlist \\
  newTrack? : TRACK
  \where
  filling = ztrue \\
  \# pool < maxPool \\
  newTrack? \notin pool \\
  owner' = owner \\
  pool' = pool \cup \{ newTrack? \} \\
  lastPlayed' = lastPlayed \\
  (mode = generatingFirst \implies \\
    \quad~ (playing' = \{ newTrack? \} \land mode' = playingFilling \\
  \quad~ \land filling' = ztrue)) \\
  (mode = playingFilling \implies \\
    \quad~ (playing' = playing \\
      \quad~ \land (\# pool' = maxPool \implies \\
      \quad\quad~ (mode' = playingRotating \land filling' = zfalse)) \\
      \quad~ \land (\# pool' < maxPool \implies \\
  \quad\quad~ (mode' = playingFilling \land filling' = ztrue))))
\end{schema}

\subsection{Advance (auto-advance and manual skip)}

When the current track ends, move to another track, avoiding an
immediate repeat of the just-ended track when the pool holds more than
one. With a single track (the transient state before the second track
lands) it replays --- the only intended looping. This one transition
serves both the automatic track-end advance and the manual
\texttt{/music next}.

\begin{schema}{Advance}
  \Delta Playlist
  \where
  (mode = playingFilling \lor mode = playingRotating) \\
  pool' = pool \\
  owner' = owner \\
  mode' = mode \\
  filling' = filling \\
  lastPlayed' = playing \\
  (\# pool \geq 2 \implies \\
    \quad~ (playing' \subseteq pool \land \# playing' = 1 \\
  \quad~ \land playing' \neq playing)) \\
  (\# pool = 1 \implies playing' = playing)
\end{schema}

The automatic track-end advance fires with no owner check. A manual
\texttt{Skip} is the same transition, additionally gated on the requester
owning music: \texttt{who?} must be the current owner. Because
\texttt{Advance} is enabled only in \texttt{playingFilling} or
\texttt{playingRotating}, both share the ``a track is really playing''
precondition -- \texttt{Skip} in \texttt{generatingFirst} is impossible
(Z finding~\#1). The implementation keys this guard off the playing track,
not the \emph{(vibe, style)} prefix glob, so a custom-named track whose stem
misses the glob still advances.

\begin{zed}
  TrackEnds \defs Advance
\end{zed}

\begin{schema}{Skip}
  Advance \\
  who? : SESSION
  \where
  who? \in owner
\end{schema}

\subsection{VibeChange}

A vibe/style change finishes the current song, then switches to the new
pool (the tracks already saved on disk for the new key). The selection
of initial mode/playing/filling mirrors \texttt{TurnOn} applied to the
new pool. \texttt{newPool?} is the new key's on-disk track set. Only the
owner may retune, and only while music is on (\texttt{mode} $\neq$
\texttt{off}) -- a forwarded vibe while off is refused, so it spends no
generation credits (finding~\#4). Ownership is unchanged.

\begin{schema}{VibeChange}
  \Delta Playlist \\
  newPool? : \power TRACK \\
  who? : SESSION
  \where
  mode \neq off \\
  who? \in owner \\
  owner' = owner \\
  \# newPool? \leq maxPool \\
  pool' = newPool? \\
  lastPlayed' = \emptyset \\
  (newPool? = \emptyset \implies \\
    \quad~ (mode' = generatingFirst \land playing' = \emptyset \\
  \quad~ \land filling' = ztrue)) \\
  (\# newPool? \geq 1 \land \# newPool? < maxPool \implies \\
    \quad~ (mode' = playingFilling \land playing' \subseteq newPool? \\
  \quad~ \land \# playing' = 1 \land filling' = ztrue)) \\
  (\# newPool? = maxPool \implies \\
    \quad~ (mode' = playingRotating \land playing' \subseteq newPool? \\
  \quad~ \land \# playing' = 1 \land filling' = zfalse))
\end{schema}

\subsection{TurnOff}

Cancel the background fill and stop playback. Ownership is released so a
later forwarded vibe from the old owner is refused (finding~\#5). The pool
remains on disk.

\begin{schema}{TurnOff}
  \Delta Playlist
  \where
  mode \neq off \\
  mode' = off \\
  playing' = \emptyset \\
  lastPlayed' = \emptyset \\
  owner' = \emptyset \\
  filling' = zfalse \\
  pool' = pool
\end{schema}

\subsection{Disable}

An unrecoverable failure (the first track could not be generated) turns
music off with the same cleanup as \texttt{TurnOff}: fill cancelled,
playback stopped, and -- crucially -- ownership released and the
avoid-repeat key (\texttt{lastPlayed}) cleared, so a disabled session
leaves nothing a later message could act on (finding~\#6). It shares
\texttt{TurnOff}'s postcondition.

\begin{zed}
  Disable \defs TurnOff
\end{zed}

\section{Key Properties}

The following hold by construction (state invariants) or as operation
postconditions; they are the intended proof obligations for animation:

\begin{itemize}
  \item \textbf{Bounded pool}: $\# pool \leq maxPool$ always.
  \item \textbf{Generation only below full}: \texttt{FillProduces} requires
    $\# pool < maxPool$; at $\# pool = maxPool$ the mode is
    \texttt{playingRotating} with $filling = zfalse$, so no generation
    occurs.
  \item \textbf{At most one fill}: $filling = ztrue$ implies
    \texttt{mode} is \texttt{generatingFirst} or \texttt{playingFilling}.
  \item \textbf{Playing is real}: when $mode \neq off$, $\# playing = 1$
    and $playing \subseteq pool$.
  \item \textbf{No immediate repeat}: \texttt{Advance} with $\# pool \geq 2$
    gives $playing' \neq playing$ (equivalently $playing' \neq
    lastPlayed'$).
  \item \textbf{Sole looping case}: \texttt{Advance} replays only when
    $\# pool = 1$, the transient before the second track lands.
  \item \textbf{Owned exactly when on}: $owner = \emptyset \iff mode =
    off$. \texttt{TurnOn} claims the owner; \texttt{TurnOff} and
    \texttt{Disable} release it, so a stopped or disabled session is
    unowned and cannot be driven by a stale message.
  \item \textbf{Control requires ownership while on}: \texttt{Skip} and
    \texttt{VibeChange} both require $who? \in owner$, and
    \texttt{VibeChange} additionally requires $mode \neq off$ -- a
    forwarded vibe while off spends no credits.
\end{itemize}

\section{Findings from Formalisation}
\label{sec:findings}

Writing the model surfaced these decisions, which the informal design
(mission m-2026-07-04-018) must resolve consistently:

\begin{enumerate}
  \item \textbf{Skip while generating the first track is undefined and is
    excluded.} \texttt{Advance} requires \texttt{playingFilling} or
    \texttt{playingRotating}; in \texttt{generatingFirst} nothing is
    playing, so \texttt{/music next} must be a no-op until track \#1
    lands.
  \item \textbf{Auto-advance and manual skip are the same transition}
    ($TrackEnds \defs Skip \defs Advance$). The implementation should
    share one code path, not two.
  \item \textbf{TurnOn, VibeChange, and restart share one ``enter a pool''
    rule} keyed on $\# pool$ (empty $\to$ generate first; partial $\to$
    play + fill; full $\to$ play + rotate). This is a natural single
    unit to extract.
  \item \textbf{VibeChange abstracts over ``finish current first.''} The
    model captures the switch after the current song ends; the temporal
    ``finish current'' is a refinement realised by deferring the switch
    to the next track-end.
  \item \textbf{The only looping is the single-track transient.} Any other
    repeat (the shipped bug) violates the \emph{No immediate repeat}
    property and is therefore a defect, not a mode.
  \item \textbf{Fresh-track precondition ties to naming.}
    \texttt{FillProduces} requires $newTrack? \notin pool$; this is only
    safe because bas7's deterministic naming guarantees uniqueness.
  \item \textbf{Ownership is a lifecycle invariant, not per-command
    state.} Modelling $owner = \emptyset \iff mode = off$ forced every
    stop path to release ownership. The shipped code cleared the owner on
    neither \texttt{TurnOff} nor \texttt{Disable}, so a stale forwarded
    vibe passed the owner check and could restart the fill (findings
    \#5/\#6). \texttt{VibeChange}'s $mode \neq off$ precondition likewise
    exposed that the code retargeted the fill even while off (finding
    \#4). Pinning ownership to the mode in the state predicate makes these
    a type-level obligation rather than scattered defensive checks.
  \item \textbf{The fill targets the played pool, not the session key.}
    Selection and fill must key on the same pool identifier. The shipped
    code selected the replayed track's pool but keyed the fill off the
    session \emph{(vibe, style)}, so a named replay grew the wrong pool
    or skipped the intended one (findings \#1/\#7). The refined
    implementation threads one pool prefix through both.
\end{enumerate}

\end{document}
